% Generated draft: requires mathematical and source review.
\documentclass[11pt]{article}
\usepackage[margin=25mm]{geometry}
\usepackage{fontspec}
\setmainfont{Latin Modern Roman}
\usepackage{amsmath,amssymb,mathtools,unicode-math}
\setmathfont{Latin Modern Math}
\usepackage{enumitem,needspace,xurl,hyperref}
\setlength{\parindent}{0pt}
\setlength{\parskip}{5pt}
\setlength{\emergencystretch}{3em}
\newcommand{\N}{\mathbb{N}}
\newcommand{\Z}{\mathbb{Z}}
\newcommand{\Q}{\mathbb{Q}}
\newcommand{\R}{\mathbb{R}}
\newcommand{\C}{\mathbb{C}}
\newcommand{\F}{\mathbb{F}}
\newcommand{\A}{\mathbb{A}}
\newcommand{\PP}{\mathbb{P}}
\newcommand{\E}{\mathbb{E}}
\newcommand{\eps}{\varepsilon}
\newcommand{\dd}{\,\mathrm{d}}
\begin{document}
\section*{UnsolvedMath ID: 30004779}
\textbf{Source problem code:} catalogue entry
\subsection{Definitions and mathematical statement}
The supplied record does not define the parameter space \(C^{\bullet}\), the morphism \(f\), the stabilizer \(P_\gamma\), or the cohomological formalism. Retain these as source-dependent data. For an algebraically closed field \(K\) and \(\gamma\in C^{\bullet}(K)\), write
\[
X_\gamma:=f^{-1}(\gamma),\qquad V_\gamma:=R\Gamma(X_\gamma,\omega_{X_\gamma}),
\]
where \(\omega_{X_\gamma}\) denotes the dualizing object in the intended, unspecified sheaf theory, and \(R\Gamma\) denotes derived global sections. Let \(\pi_0(P_\gamma)\) denote the group of connected components of the stabilizer, and let \(q_\gamma:P_\gamma\to\pi_0(P_\gamma)\) be the component projection, whenever these are defined in the intended setting.

\textbf{Conjecture (source formulation).} For every algebraically closed field \(K\) and every \(\gamma\in C^{\bullet}(K)\), the specified action of \(P_\gamma\) on \(V_\gamma\) factors through \(q_\gamma\): there exists an action of \(\pi_0(P_\gamma)\) on \(V_\gamma\) whose pullback along \(q_\gamma\) is the given action.

This concerns the action on the derived object \(V_\gamma\), rather than merely actions on its individual cohomology groups. The source record supplies no characteristic restrictions on \(K\). The spaces, group objects, coefficients, existence of the dualizing object, and precise meaning of an action and its factorization require clarification before this becomes a fully specified mathematical statement.
\subsection{Short English statement}
The conjecture asks whether, for every algebraically closed field and every allowed parameter, the stabilizer action on the derived global sections of the fiber's dualizing object depends only on the stabilizer's connected component. The underlying geometric and cohomological setting is missing from the supplied record.
\subsection{Sources}
\begin{itemize}
\item Local UnsolvedMath registry, permanent ID \texttt{30004779}, title \emph{Component-Group Factorization on Dualizing Cohomology}, source problem code \texttt{catalogue entry}, in the Oberwolfach Reports Open Problems collection. The record gives proposed year 2021 and supplies no individual proposer.
\item Supplied report citation: \emph{Automorphic Forms, Geometry and Arithmetic (2022)}, Oberwolfach Reports, pp.\ 2089--2156, DOI \texttt{10.4171/owr/2021/39}. The literature-source label is \emph{Oberwolfach Report 39/2021}. \url{https://doi.org/10.4171/owr/2021/39}.
\item Supplied literature reference: Bouthier, \emph{Faisceaux caracteres}. \url{https://arxiv.org/abs/2108.12663}.
\item Literature assessment dated exactly \texttt{2026-08-22}: the supplied curation labels the problem open, reports that this factorization conjecture is an assumption for a perversity theorem in the report, and states that no general proof was found. This is supplied dated curation, not an independently verified literature assessment. The background explicitly requests independent verification of time-sensitive status and cited claims.
\end{itemize}
\end{document}
